\subsection{Retention under discrete stochastic optimizer updates}
\label{app:theory-optimizer}

For discrete noisy updates, the useful quantity is the total upward drift
allowed in the joint loss. Controlling that budget permits individual steps
to increase the loss while keeping a high-probability ceiling on the loss
over a stated update horizon.

Lemma~\ref{lem:shared-exemplar-retention} converts a deterministic energy
ceiling into complete-exemplar protection. We next replace pathwise descent
with a controlled conditional drift, allowing shared parameters, noise,
bias, and predictable preconditioning. The descent estimate follows the
variable-metric analysis of \citet[Section~2]{wang2017sqn}; the trajectory
bound uses Ville's nonnegative-supermartingale inequality, as stated by
\citet[Lemma~1]{howard2020timeuniform}. These are standard ingredients;
the result below applies them to the verified-response potential.

\begin{theorem}[Conditional stochastic exemplar retention]
\label{ext:thm-stochastic-retention}
Fix finitely many complete verified responses $(x_j,e_j)$ realizing keys
$b_j$, with lengths $L_j>0$ and weights $\alpha_j>0$. Define
\[
 R(\theta)=\sum_j\frac{\alpha_j}{L_j}
             [-\log\pi_\theta(e_j\mid x_j)],\qquad
 F(\theta)=R(\theta)-J(\theta),\qquad J(\theta)\le J_{\max}<\infty.
\]
Assume $F$ is differentiable with globally $L_F$-Lipschitz gradient and
$\theta_0$ is deterministic with finite $F(\theta_0)$. Relative to the
history filtration $(\mathcal F_t)$, consider
\[
 \theta_{t+1}=\theta_t-\eta_t H_t(g_t+b_t+\xi_t),
 \qquad g_t=\nabla F(\theta_t).
\]
The variables $\theta_t,\eta_t,H_t,b_t$ are $\mathcal F_t$-measurable,
$0\preceq H_t\preceq MI$, $\eta_t>0$, and $\eta_t L_FM\le1$.
The noise satisfies $\mathbb E[\xi_t\mid\mathcal F_t]=0$ and has finite
conditional second moment. Suppose deterministic $d_t\ge0$ satisfy
\begin{equation}
 \frac{\eta_t}{2}\|b_t\|_{H_t}^{2}
 +\frac{L_F\eta_t^2}{2}
       \mathbb E[\|H_t\xi_t\|_2^2\mid\mathcal F_t]\le d_t,
 \qquad \|u\|_H^2=u^\top Hu.
 \label{ext:eq-optimizer-budget}
\end{equation}
For $N\in\mathbb N\cup\{\infty\}$ with
$B_N=\sum_{t<N}d_t<\infty$, write $D_0=F(\theta_0)+J_{\max}\ge0$.
For every $0<\delta<1$, with probability at least $1-\delta$, the following
holds simultaneously for all integer $t\le N$ and all protected exemplars:
\begin{equation}
 p_{\theta_t}(b_j\mid x_j)\ge\pi_{\theta_t}(e_j\mid x_j)
 \ge\exp\!\left[-\frac{L_j(D_0+B_N)}{\alpha_j\delta}\right].
 \label{ext:eq-stochastic-retention}
\end{equation}
For $N=\infty$, each protected probability also has an almost surely
positive, potentially random, infimum over all iterations. The likelihood
must include the complete response and its termination event, or a fixed
deterministic horizon, under the sampling law whose key probability is bounded.
\end{theorem}
\begin{proof}
Smoothness, conditional centering of $\xi_t$, and
$H_t^2\preceq MH_t$ imply
\begin{align*}
 \mathbb E[F(\theta_{t+1})\mid\mathcal F_t]
 &\le F(\theta_t)-\eta_tg_t^\top H_t(g_t+b_t)
       +\frac{\eta_t}{2}\|g_t+b_t\|_{H_t}^2\\
 &\hspace{1.5em}+\frac{L_F\eta_t^2}{2}
          \mathbb E[\|H_t\xi_t\|_2^2\mid\mathcal F_t]\\
 &\le F(\theta_t)-\frac{\eta_t}{2}\|g_t\|_{H_t}^2+d_t.
\end{align*}
The bias cross term cancels when the square is expanded. Since
$F+J_{\max}\ge R\ge0$, the process
\[
 X_t=F(\theta_t)+J_{\max}+\sum_{s=t}^{N-1}d_s
\]
is a nonnegative supermartingale, using the convergent infinite tail when
$N=\infty$. Its initial value is $D_0+B_N$. For $D_0+B_N>0$, Ville's
inequality gives
$\Pr(\sup_{t\le N}X_t\ge(D_0+B_N)/\delta)\le\delta$.
If $D_0+B_N=0$, nonnegativity gives $X_t=0$ almost surely instead.
Every nonnegative summand of $R(\theta_t)\le X_t$ obeys the same ceiling.
Exponentiation therefore yields~\eqref{ext:eq-stochastic-retention}
without a union factor over exemplars. A complete-exemplar event is contained
in its verified-key event. For $N=\infty$, the same maximal bound gives
$\Pr(\sup_tX_t\ge u)\le(D_0+B_\infty)/u$ for $u>0$.
Letting $u\to\infty$ proves $\sup_tX_t<\infty$ almost surely and hence
random positive all-time floors.
\end{proof}

\paragraph{Finite and infinite horizons.}
With zero bias and uniformly bounded noise variance, the cost in
\eqref{ext:eq-optimizer-budget} is $O(\eta_t^2)$, so square-summable steps
suffice for an infinite-horizon budget. Bias instead contributes
$O(\eta_t\|b_t\|^2)$. Constant steps with persistent variance yield the
finite-horizon result, with a budget that can grow with $N$; this argument
does not give infinite-horizon protection in that regime. The classical
almost-supermartingale framework of \citet{robbins1971almost} supplies
related convergence tools, but the proof above needs only the displayed
drift and maximal inequality. No positive lower eigenvalue of $H_t$ is
needed for retention alone, and the theorem does not assert stationarity,
parameter convergence, or a useful evaluation-scale probability floor.

\paragraph{Example: a noise budget for one protected response.}
Consider a hypothetical fixed one-token horizon with two possible complete
responses and
$\pi_\theta(e)=(1+e^{-\theta})^{-1}$, $L=\alpha=1$, and
$F=R=-\log\pi_\theta(e)$, with $J=J_{\max}=0$. Start at $\theta_0=0$, so
$D_0=\log2$. This scalar objective has $L_F=1/4$ and $|g_t|\le1$.
Take $H_t=1$, zero bias, and independent noise $\xi_t\in\{-0.1,0.1\}$ with
equal probabilities. At constant step $\eta_t=0.1$, the theorem permits
$d_t=1/80000$, giving $B_{100}=0.00125$ but $B_\infty=\infty$.
With $\eta_t=0.1/(t+1)$ instead,
$B_\infty=\pi^2/480000\approx2.06\times10^{-5}$, which supplies the
infinite-horizon condition. For the constant-step, 100-update case, the
$95\%$ probability floor from~\eqref{ext:eq-stochastic-retention} is only
$\exp[-(\log2+0.00125)/0.05]\approx9.30\times10^{-7}$.
The next result uses bounded noise to improve this confidence dependence.

\begin{corollary}[Sharper confidence dependence under bounded noise]
\label{ext:cor-bounded-noise-retention}
Under the assumptions of Theorem~\ref{ext:thm-stochastic-retention},
suppose $b_t=0$, $N<\infty$,
and $\eta_t$ is deterministic. Let deterministic bounds $G_t,s_t<\infty$
satisfy $\|g_t\|_2\le G_t$ and $\|\xi_t\|_2\le s_t$ almost surely. Define
\[
 c_t=\eta_tMG_ts_t,\quad V_N=\sum_{t<N}c_t^2,\quad
 Q_N=\frac{L_FM^2}{2}\sum_{t<N}\eta_t^2s_t^2.
\]
For every $0<\delta<1$, with probability at least $1-\delta$,
simultaneously for all $t\le N$ and $j$,
\begin{equation}
 p_{\theta_t}(b_j\mid x_j)\ge\pi_{\theta_t}(e_j\mid x_j)
 \ge\exp\!\left[-\frac{L_j}{\alpha_j}
   \left(D_0+Q_N+\sqrt{2V_N\log(1/\delta)}\right)\right].
 \label{ext:eq-bounded-noise-retention}
\end{equation}
\end{corollary}
\begin{proof}
The pathwise smoothness estimate gives
\[
 F(\theta_{t+1})-F(\theta_t)
 \le-\frac{\eta_t}{2}\|g_t\|_{H_t}^2+Z_{t+1}
       +\frac{L_F\eta_t^2M^2s_t^2}{2},
\]
where
\[
 Z_{t+1}=-\eta_t g_t^\top H_t(I-L_F\eta_tH_t)\xi_t.
\]
Predictability makes $Z_{t+1}$ a martingale difference. The spectral bounds
imply $0\preceq I-L_F\eta_tH_t\preceq I$, so $|Z_{t+1}|\le c_t$.
Conditional Hoeffding's lemma therefore makes
\[
 \exp\!\left(\lambda\sum_{s<t}Z_{s+1}
              -\frac{\lambda^2}{2}\sum_{s<t}c_s^2\right),\qquad\lambda>0,
\]
a nonnegative supermartingale starting at one. This standard exponential
construction is reviewed by \citet{howard2020timeuniform}. Ville's inequality
and optimization over $\lambda$ yield
\[
 \Pr\!\left(\max_{t\le N}\sum_{s<t}Z_{s+1}
              >\sqrt{2V_N\log(1/\delta)}\right)\le\delta
\]
when $V_N>0$; when $V_N=0$, the sum is identically zero. Sum the drift inequalities, drop the nonpositive gradient terms,
and apply the summand bound from Theorem~\ref{ext:thm-stochastic-retention}.
\end{proof}

In the same 100-update example, $G_t=M=1$ and $s_t=0.1$ give
$V_{100}=0.01$ and $Q_{100}=0.00125$. The sharper $95\%$ floor is therefore
\[
 \exp\!\left[-\log2-0.00125-\sqrt{0.02\log20}\right]\approx0.391.
\]
Both bounds apply simultaneously to the toy trajectory's first 100 updates;
the improvement comes from stronger noise assumptions. These toy constants are assigned for the calculation; they do not
estimate noise or curvature in the language-model runs.

\paragraph{Relation to stochastic policy-gradient guarantees.}
\citet{zhang2021reinforce} analyze stochastic tabular REINFORCE with a
state-action log barrier, phased step and regularization schedules, and
policy mixing. Their regret guarantees provide a related stochastic
policy-gradient precedent; they do not provide a fixed per-response floor
for our replay bank. For the implemented optimizer, a current-gradient Adam
second moment is not predictable relative to that gradient's noise.
Momentum, clipping, prompt scheduling, and changing banks likewise require
separate bias or energy control. The counterexamples of
\citet{reddi2018adam} explain why an unrestricted Adam guarantee cannot be
assumed; they do not establish failure of these runs. The hypotheses above
have not been established for the experiments' AdamW/PPO updates. In
particular, positive replay weight and a small nominal learning rate do not
verify the smoothness, predictability, or accumulated-error conditions.

